%=====================================================================
% FRONT MATTER
%=====================================================================
\pagenumbering{roman}
\structuralfigures{F.}
\thispagestyle{empty}

\begin{center}
\vspace*{1.0cm}

{\color{secondary}\rule{0.85\textwidth}{2pt}}\\[6pt]
{\large\textsc{\color{secondary}Bachelor of Science in Information Technology}}\\[18pt]

{\Huge\bfseries\color{primary}DATABASE SYSTEMS}\\[10pt]
{\LARGE\bfseries\color{primary}Designing, Querying and Trusting\\[4pt]a Relational Database}\\[16pt]

{\color{secondary}\rule{0.85\textwidth}{2pt}}\\[16pt]

{\Large\textsc{Unified Lecture Handout}}\\[6pt]
{\large\color{gray!70!black}Complete Course in One Volume $\bullet$ Fall 2026}\\[24pt]

% --- the claim the whole book rests on ---
\begin{tikzpicture}
  \node[rounded corners=4pt, draw=primary, line width=1.1pt, fill=primary!6,
        text=primary, align=center, font=\small\bfseries,
        minimum width=11.6cm, minimum height=1.5cm]
       {\faIcon{database}\\[2pt]A database is not a place to put data.
        It is a set of promises about data.};
\end{tikzpicture}\\[10pt]
{\small\itshape\color{gray!70!black}This course is how each promise is made
--- and, every time, what the database does when you do not ask it to keep
one.}

\vspace{1.0cm}

\begin{tcolorbox}[enhanced, width=0.86\textwidth, colback=lightbg,
  colframe=primary, boxrule=0pt, leftrule=4pt, arc=2pt, top=6pt, bottom=6pt]
\small
\textbf{\color{primary}What this volume covers.} Every clause of the HEC
outline for \emph{Database Systems}, traced chapter by chapter in the Course
Specification overleaf.

\smallskip
That outline is published in the \textbf{2017} curriculum (p.~88) and
reprinted \emph{word for word} in the \textbf{2023} curriculum (p.~58) --- the
two differ by a single added full stop across ninety words, six years apart.
The \textbf{2025} curriculum deletes it: it covers only BS Computer Science,
publishes no course content for any major course at all, and supplies five
rewritten learning outcomes with no Bloom levels. All of that is recorded in
the Course Specification rather than smoothed over.

\smallskip
The course is taught on a \textbf{running database}. One schema is designed in
Chapter~6, created in Chapter~10, and every chapter afterwards adds to it or
asks something of it; by Chapter~23 it is the semester project. The tool is
\textbf{SQLite}, which you already have.
\end{tcolorbox}

\vspace{0.9cm}
{\small\color{gray!60!black}Six Parts $\bullet$ Twenty-Three Chapters $\bullet$ Seven Appendices}
\vspace{1cm}
\end{center}

\newpage

%=====================================================================
\chapter*{How to Use This Handout}
\addcontentsline{toc}{chapter}{How to Use This Handout}
\markboth{How to Use This Handout}{}

There are two ways to teach this subject badly.

The first is to teach it as theory. Relational algebra on a whiteboard,
normal forms as definitions to be recited, ACID as four words beginning with
the right letters. A student who has learned it this way can answer an
examination question about third normal form and cannot look at a table and
say whether it is in it.

The second is to teach it as a tool. Open the software, click here, type a
\sqlkw{SELECT}, admire the rows. A student who has learned it this way can
produce an answer and has no way of knowing whether it is the right one,
because nothing in the course ever told them what the database was promising.

\textbf{This handout does neither, and the reason is in the title.} A database
makes promises --- that this attribute is a number, that this row refers to a
row that exists, that either all of this happened or none of it did --- and
every one of those promises is \emph{checkable}. So the method of this book is
to state the promise, then ask the database to break it, and see what happens.
Eleven chapters end with a figure produced that way on a real machine. They
are marked, and the file that produces each one is named.

\section*{The furniture}

Every chapter has the same parts, in the same order, so that you can find
what you need without reading from the beginning.

\rowcolors{2}{white}{lightbg}
\begin{longtable}{@{}L{4.2cm}L{10.4cm}@{}}
\toprule
\rowcolor{primary!15}
\textbf{You will see} & \textbf{It means} \\
\midrule
\endhead
\faIcon{bullseye}~Learning Outcomes & What you should be able to \emph{do}
afterwards. Five per chapter, each with a verb you can be examined on \\
\faIcon{link}~Before You Start & What this chapter assumes, and where it was
taught \\
\faIcon{tags}~Key Terms & Every term defined here. If one is unfamiliar at the
end, the chapter did not work \\
\faIcon{book}~Definition & A term being given its precise meaning. These are
the sentences to learn exactly \\
\faIcon{lightbulb}~Concept & An idea, a rule of thumb, or a way of deciding \\
\faIcon{triangle-exclamation}~Watch Out & Something that is true, surprising,
and will cost you if you do not know it \\
\faIcon{code}~Example & A worked case, usually from the running database \\
\faIcon{pen-to-square}~Worked & Arithmetic, every step shown. Nothing is left
``by inspection'' \\
\faIcon{tachometer-alt}~Measured & A figure produced by running the named file
in \texttt{code/}, on the machine in Appendix~A. You can reproduce every one \\
\faIcon{circle-exclamation}~Common Pitfalls & What goes wrong, collected, so
that you meet it here rather than in the laboratory \\
\faIcon{flask}~Try It Yourself & The laboratory. Numbered steps, run in order \\
\faIcon{circle-question}~Checkpoint & Three questions. Answers in Appendix~C,
and they are full answers rather than a line each \\
\faIcon{list-check}~What to Take Away & The chapter in eight or nine lines \\
\faIcon{pen-fancy}~Review Questions & Seven: two multiple choice, three short,
two applied \\
\bottomrule
\end{longtable}

\section*{The laboratory, and what it needs}

\textbf{Nothing to install.} The whole course runs on \textbf{SQLite}, which
is already on every machine you are likely to use: it ships inside Python, it
ships with most operating systems, and a database is a single file you can
copy, email or delete. Appendix~A sets it up in one sitting, which for most
readers means checking that it is already there.

That choice is not a compromise with the curriculum. \textbf{HEC names no
database product in any edition} --- the 2017 and 2023 outlines say only
``in any DBMS'', and 2025 says ``popular database management systems''. What
SQLite can and cannot show you is stated plainly in \chref{whynotfiles} and
again in \chref{project}, because a tool's limits are part of what you are
learning.

\section*{The database we build}

One schema, carried through the whole course. You design it in
\chref{ermodel} and \chref{cardinality}, map it to tables in \chref{mapping},
create it in \chref{ddl}, fill it, query it for five chapters, find out in
\chref{normalforms} that it was wrong, fix it, index it, and defend it. It is
a university enrolment database --- students, courses, sections, instructors,
enrolments, grades --- chosen because you already know how it is supposed to
behave, so no lecture time goes on explaining the subject matter.

\section*{If you read nothing else}

\begin{tight}
  \item \chref{constraints}, on the difference between \emph{declaring} a key
        and \emph{having} one.
  \item \chref{normalforms}, which is the chapter this course exists for.
  \item \chref{indexes}, where the rule everybody repeats stops being true and
        you can see the exact row count at which it does.
\end{tight}

\newpage

%=====================================================================
\chapter*{Course Specification}
\addcontentsline{toc}{chapter}{Course Specification}
\markboth{Course Specification}{}

This chapter exists so that a student, a lecturer or a reviewer can see
exactly what this handout claims to cover and on what authority. Everything in
it is traceable to a published document, and where those documents are wrong,
unclear or contradict one another, that is recorded here rather than quietly
corrected.

\section*{Particulars}

\rowcolors{2}{white}{lightbg}
\begin{longtable}{@{}L{3.6cm}L{11.0cm}@{}}
\toprule
\rowcolor{primary!15}
\textbf{Item} & \textbf{Value} \\
\midrule
\endhead
Course title & \emph{Database Systems}. HEC 2023's description page calls it
\emph{Database System}, singular, while every table in the same booklet says
\emph{Systems}. The plural is used here \\
Course code & \textbf{CS-103}. Unlike every sibling course in this series,
there is no \texttt{INFT-} equivalent: the registrar uses CS-103 and nothing
else \\
Programme & BS Information Technology, \textbf{2nd semester}. This matches HEC
2023, which places the course in semester~2 of every computing degree. HEC
2017 placed it at BSIT semester~V; the discrepancy is recorded, not resolved \\
Credit hours & \textbf{4 (3+1)} --- three lecture credits and \textbf{one
credited laboratory}. The laboratory is paid for, which is why every chapter
has one and why they are long enough to be worth the hour \\
Prerequisite & \textbf{None.} HEC 2017's own booklet gives two different
answers --- see defect~4 \\
Duration & Sixteen teaching weeks, midterm in Week~8, project due Week~16
(Appendix~G) \\
Teaching methodology & As published (HEC 2017): lectures, written assignments,
\textbf{practical labs}, semester project, presentations \\
Assessment & As published (HEC 2017): sessional exam, home assignments,
quizzes, project, presentations, final exam \\
Set texts & As published: Connolly and Begg; Garcia-Molina, Ullman and Widom;
Silberschatz, Korth and Sudarshan; Ramakrishnan and Gehrke. No publication
years are given for any of the four, in any edition --- see Appendix~E \\
\bottomrule
\end{longtable}

\section*{Where the outline comes from}

\rowcolors{2}{white}{lightbg}
\begin{longtable}{@{}L{1.6cm}L{4.4cm}L{3.2cm}L{5.0cm}@{}}
\toprule
\rowcolor{primary!15}
\textbf{Edition} & \textbf{Status} & \textbf{Semester} &
\textbf{Course content published?} \\
\midrule
\endhead
\textbf{2017} & Computing Core --- compulsory in BSCS, BSSE \emph{and} BSIT &
IV, IV, \textbf{V} & \textbf{Yes} --- p.~88 \\
\textbf{2023} & Computing Core --- compulsory in all ten computing degrees and
the Associate Degree & \textbf{2} & \textbf{Yes} --- p.~58, \emph{the 2017 text
reprinted} \\
2025 & Major Compulsory --- but the edition covers only BS Computer Science
and the Associate Degree. \textbf{There is no BSIT in it at all} & III &
\textbf{No} --- deleted; five CLO bullets only \\
\bottomrule
\end{longtable}

\begin{alertbox}[Six defects in the published documents]
Recorded rather than inherited. Each is stated with the correction this
handout uses.

\begin{tightnum}
  \item \textbf{The syllabus is frozen.} The 2023 course outline is the 2017
        course outline word for word, with \textbf{one added full stop}
        (``vs'' becomes ``vs.'') across ninety words --- six years apart. Then
        2025 removes it entirely. A field that changed a great deal in that
        period is described by a text that did not change at all.
  \item \textbf{The outline normalises before it designs.} The clause list
        reaches \emph{normalization, functional dependencies, normal forms}
        \textbf{before} \emph{entity relationship model \dots\
        entity-relationship diagrams}. You cannot normalise a schema you have
        not designed. All four of HEC's own reference books teach the
        other order. \textbf{This handout reorders}: design in Part~II, map in
        Part~III, normalise in Part~V. The clause numbers in the table below
        show exactly where each one moved to.
  \item \textbf{One clause is printed twice.} ``attributes'' appears in the
        relational-model run and again in the entity-relationship run. Thirty
        three clauses, thirty two distinct. Both occurrences are traced below
        and the second is marked as a duplicate.
  \item \textbf{The prerequisite contradicts itself inside one booklet.} HEC
        2017's course description page says \texttt{Prerequisites: None}. All
        three of its scheme-of-studies tables --- BSCS p.~25, BSIT p.~31, BSSE
        p.~39 --- say \textbf{Data Structures \& Algorithms}. At 2nd semester
        the department cannot assume Data Structures, so \textbf{this handout
        assumes none}, and \chref{indexes} and \chref{transactions} are written
        to be self-contained because of it.
  \item \textbf{The 2023 learning-outcome table is broken in print.} CLO-3's
        text stops mid-sentence; its remainder (``anomalies by normalizing
        database tables.'') is printed as CLO-4's text; and the real CLO-4 is
        left in a trailing row with \emph{no number and no Bloom level}. This
        is a typesetting fault in the published national curriculum, not a
        misreading. The outcomes are reconstructed below from the 2017 text,
        which is intact.
  \item \textbf{A Bloom level is wrong.} CLO-3 reads ``Identify functional
        dependencies and \emph{resolve} database anomalies by normalizing
        database tables'' and is tagged \textbf{C2}. Identifying is C2;
        resolving and normalising are C3 at least. Every sibling course in
        this series complains that HEC left the Bloom column blank. Here it is
        filled in --- and filled in incorrectly, which is harder to notice.
\end{tightnum}
\end{alertbox}

\section*{Course content, traced clause by clause}

The outline is reproduced below exactly as published, in its published order,
split into its thirty-three clauses. The right-hand column is where each is
taught. Clauses that move are the subject of defect~2.

\rowcolors{2}{white}{lightbg}
\begin{longtable}{@{}L{0.9cm}L{9.3cm}L{4.4cm}@{}}
\toprule
\rowcolor{primary!15}
\textbf{\#} & \textbf{Clause, as published} & \textbf{Taught in} \\
\midrule
\endhead
1 & Basic database concepts & \chref{whynotfiles} \\
2 & Database approach vs file based system & \chref{whynotfiles} \\
3 & database architecture & \chref{whodoeswhat} \\
4 & three level schema architecture & \chref{architecture} \\
5 & data independence & \chref{architecture} \\
6 & relational data model & \chref{relational} \\
7 & attributes & \chref{relational} \\
8 & schemas & \chref{relational} \\
9 & tuples & \chref{relational} \\
10 & domains & \chref{relational} \\
11 & relation instances & \chref{relational} \\
12 & keys of relations & \chref{constraints} \\
13 & integrity constraints & \chref{constraints} \\
14 & relational algebra & \chref{algebra} \\
15 & selection & \chref{algebra} \\
16 & projection & \chref{algebra} \\
17 & Cartesian product & \chref{algebra} \\
18 & types of joins & \chref{joins} \\
19 & \textbf{normalization} & \chref{normalforms} \emph{(moved later)} \\
20 & \textbf{functional dependencies} & \chref{fds} \emph{(moved later)} \\
21 & \textbf{normal forms} & \chref{normalforms} \emph{(moved later)} \\
22 & \textbf{entity relationship model} & \chref{ermodel} \emph{(moved
earlier)} \\
23 & entity sets & \chref{ermodel} \\
24 & attributes & \textbf{duplicate of 7} \\
25 & relationship & \chref{ermodel}, \chref{cardinality} \\
26 & entity-relationship diagrams & \chref{cardinality}, \chref{mapping} \\
27 & Structured Query Language (SQL) & \chref{ddl}, \chref{select} \\
28 & Joins and sub-queries in SQL & \chref{joins}, \chref{subqueries} \\
29 & Grouping and aggregation in SQL & \chref{grouping} \\
30 & concurrency control & \chref{transactions} \\
31 & database backup and recovery & \chref{transactions} \\
32 & indexes & \chref{indexes} \\
33 & NoSQL systems & \chref{project} \\
\bottomrule
\end{longtable}

\begin{conceptbox}[Three topics taught that the outline does not name]
A reviewer will ask, so the justification is here rather than implied.

\begin{tight}
  \item \textbf{\chref{evolution}, changing a schema that is already full.}
        The outline never mentions \sqlkw{ALTER}. A student who cannot change
        a schema cannot carry one database through sixteen weeks, and the
        project requires exactly that.
  \item \textbf{\chref{views}, views and triggers.} Absent from the clause
        list. Required by CLO-2, which asks for \emph{logical and physical}
        schemas, and assumed by the department's own
        \emph{E-Commerce Application Development} (INFT-4041), which lists
        this course as its prerequisite --- see Appendix~F.
  \item \textbf{\chref{project}, security and authorisation.} Absent from 2017
        and 2023 entirely. \textbf{HEC 2025 requires it by name}: its fifth
        learning outcome is ``Analyze the role of databases in information
        systems and their security considerations''. The newest edition
        justifies the addition.
\end{tight}
\end{conceptbox}

\section*{Course Learning Outcomes}

\textbf{HEC 2017's four, verbatim}, with the Domain and Bloom levels exactly as
published --- including the one that is wrong.

\rowcolors{2}{white}{lightbg}
\begin{longtable}{@{}L{1.4cm}L{9.6cm}L{1.6cm}L{2.0cm}@{}}
\toprule
\rowcolor{primary!15}
\textbf{CLO} & \textbf{As published} & \textbf{Level} & \textbf{Taught in} \\
\midrule
\endhead
CLO-1 & ``Explain fundamental database concepts.'' & C2 & Part~I \\
CLO-2 & ``Design conceptual, logical and physical database schemas using
different data models.'' & C5 & Parts~II, III \\
CLO-3 & ``Identify functional dependencies and resolve database anomalies by
normalizing database tables.'' & \textbf{C2}\,$^\dagger$ & Part~V \\
CLO-4 & ``Use Structured Query Language (SQL) for database definition and
manipulation in any DBMS.'' & C4 & Parts~III, IV \\
\bottomrule
\end{longtable}

\noindent$^\dagger$\,\textbf{As published, and incorrect} --- see defect~6.
Assessed in this course at \textbf{C3}.

\medskip
\noindent\textbf{Four added by the department}, each discharging one of HEC
2025's five outcomes, which the 2017 clause list does not reach.

\rowcolors{2}{white}{lightbg}
\begin{longtable}{@{}L{1.4cm}L{1.3cm}L{8.3cm}L{3.6cm}@{}}
\toprule
\rowcolor{primary!15}
\textbf{CLO} & \textbf{Level} & \textbf{Outcome} & \textbf{Discharges} \\
\midrule
\endhead
CLO-5 & C3 & Write SQL for retrieval, insertion, update and deletion against a
schema you did not design & HEC 2025 outcome~2 \\
CLO-6 & C2 & Explain transactions, concurrency and recovery, and demonstrate
what a transaction prevents & HEC 2025 outcome~3 \\
CLO-7 & C3 & Implement and administer a working database, including backup and
restore & HEC 2025 outcome~4 \\
CLO-8 & C4 & Analyse a database's role in a wider system and its security
posture, and justify the analysis from evidence & HEC 2025 outcome~5 \\
\bottomrule
\end{longtable}

\newpage

%=====================================================================
\chapter*{Prerequisite Self-Check}
\addcontentsline{toc}{chapter}{Prerequisite Self-Check}
\markboth{Prerequisite Self-Check}{}

This is a \textbf{second-semester} course. It assumes one semester of
programming and nothing else --- in particular it does \emph{not} assume Data
Structures, whatever HEC 2017's scheme-of-studies tables say (defect~4).

Twelve questions. If you can answer ten, start at \chref{whynotfiles}. If you
cannot, the right-hand column says where to look --- and none of it will take
more than an evening.

\rowcolors{2}{white}{lightbg}
\begin{longtable}{@{}L{0.8cm}L{9.4cm}L{4.4cm}@{}}
\toprule
\rowcolor{primary!15}
\textbf{\#} & \textbf{Question} & \textbf{If not} \\
\midrule
\endhead
1 & What is the difference between a file and a folder? & --- \\
2 & A text file holds one record per line, fields separated by commas. How
would you find every line whose third field is ``2026''? & Programming
Fundamentals \\
3 & Write a loop that reads a file line by line. & Programming Fundamentals \\
4 & What does it mean for two values to be \emph{equal} when one is the number
7 and the other is the text ``7''? & \chref{ddl} returns to this \\
5 & What is an array index? What happens if you use one that is too large? &
Programming Fundamentals \\
6 & Given sets $A=\{1,2,3\}$ and $B=\{2,3,4\}$, write $A \cup B$, $A \cap B$
and $A - B$. & Appendix~B \\
7 & What does $A \times B$ mean for two sets, and how many elements does it
have? & Appendix~B \\
8 & Is the statement ``$x > 5$ and $x < 3$'' ever true? & Appendix~B \\
9 & What is the truth value of ``not (P and Q)'' when P is true and Q is
false? & Appendix~B \\
10 & Two people edit the same file at the same time and both save. What
happens to the first person's changes? & \chref{transactions} \\
11 & Roughly how many times do you have to halve 1{,}000{,}000 to get to 1? &
Appendix~B; \chref{indexes} needs it \\
12 & What is a function, in the mathematical sense --- how many outputs may one
input have? & Appendix~B; \chref{fds} depends on it \\
\bottomrule
\end{longtable}

\begin{conceptbox}[Question 12 is the one that matters]
Everything in Part~V rests on it. A \emph{functional dependency} is nothing
more than the statement that one set of attributes is a function of another:
give me a student's roll number and there is exactly one name it can be paired
with. If question~12 was uncomfortable, read Appendix~B before \chref{fds},
not during it.
\end{conceptbox}
